\chapterwithopening{Methodology}{chap:methodology}{
  \chaptercontentsitem{sec:methodology_introduction}
  \chaptercontentsitem{sec:system_development_methodology}
  \chaptercontentsitem{sec:development_workflow_structure}
  \chaptercontentsitem{sec:chapter3_summary}
}

\section{Introduction}\label{sec:methodology_introduction}

Choosing a development methodology for EduVerse was not a routine decision.
The platform's modules, including organization governance, class operation,
learning content, assessment, communication, AI support, and integrated IDE
support, are not independent; each depends on shared permissions, shared
data, and the roles of the Organization Administrator, Teacher, and Student.
A methodology that fixed requirements and structure at the outset, before
these dependencies were understood, risked producing modules that worked in
isolation but conflicted once integrated.

This chapter presents the methodological reasoning behind EduVerse's
development: what approach was chosen, why it suited a multi-organization
educational platform, and how it was applied consistently across features.
Engineering practices and tool-specific execution are addressed separately
in Chapter~6.

\section{System Development Methodology}\label{sec:system_development_methodology}

EduVerse was developed using a hybrid iterative methodology with
predominantly feature-based vertical-slice development, rather than a
strict waterfall model in which analysis, design, implementation, and
validation each occur once in a fixed sequence. The platform instead evolved
through repeated increments, each centered on a meaningful academic or
administrative capability, allowing requirements and structure to be refined
as the relationships between modules became clearer through development.

\subsection{Feature-Based Iterative Methodology}\label{subsec:feature_based_iterative_workflow}

Development was organized around functional areas rather than technical
layers. Iterations typically focused on modules such as organization
management, class management, materials, assignments, examinations, live
sessions, results, AI Agent support, integrated IDE support, feature
configuration, and public access workflows. Within each iteration, the
selected feature moved through recurring stages of analysis, design,
development, and validation, with the outcome of one increment informing the
next.

This structure also governed coordination between the web application, which
served as the primary client, and the MVP mobile companion application.
Because the mobile app targeted selected high-frequency academic actions
rather than full feature parity, each iteration served as a checkpoint for
deciding which workflows required cross-client refinement and which remained
web-centered.

\subsection{Vertical-Slice Development Approach}\label{subsec:vertical_slice_development_approach}

A feature was generally treated as a coherent functional unit rather than
split into isolated interface-only or backend-only stages. Depending on its
scope, a slice could span user-interface behavior, business rules, backend
or API behavior, permission handling, and data requirements. This suited a
role-based platform: an action such as submitting an assignment or
scheduling an examination cannot be understood as a user-interface concern
alone, since it also depends on who may perform it, what rules govern it,
and how the resulting record is stored and retrieved.

This description is deliberately framed as predominantly vertical-slice
rather than absolute. Some features required only a subset of system
layers, while others extended across several and, in selected cases, into
the MVP mobile companion application as well.

\subsection{Justification for Methodology Selection}\label{subsec:justification_for_methodology_selection}

A rigid sequential methodology was a poor fit for EduVerse because the
platform's module boundaries and dependencies were not fully knowable in
advance; they became clear through implementation. Organization governance,
class operation, learning content, assessment, communication, AI support,
and integrated IDE support influence one another through shared permissions,
shared data, and shared user context, so fixing their design before that
interplay was understood would have risked structural rework later.

The iterative, vertical-slice approach let the project refine permissions,
workflows, interface behavior, and data structures incrementally, checking
after each increment whether a feature remained coherent with the wider
platform before expanding further. This mattered in particular because the
Organization Administrator, Teacher, and Student each interact with related
features from different perspectives, so role-based access and workflow
consistency needed repeated, not one-time, validation.

\section{Iterative Development Structure}\label{sec:development_workflow_structure}

Although development was iterative, each increment followed the same
underlying cycle, adjusted in scope and depth to the complexity of the
feature involved. Some increments centered mainly on the web application;
others required supporting changes to the backend, database, permissions, or
MVP mobile companion application. What stayed constant was not the shape of
each iteration but the sequence it followed, moving from clarification,
through realization and validation, to integration.

\begin{table}[htbp]
  \centering
  \small
  \renewcommand{\arraystretch}{1.15}
  \setlength{\tabcolsep}{5pt}
  \caption{Repeated Iterative Development Structure in EduVerse}
  \label{tab:eduverse_iterative_development_structure}
  \begin{tabular}{|>{\raggedright\arraybackslash}p{0.27\textwidth}|>{\raggedright\arraybackslash}p{0.57\textwidth}|}
    \hline
    \textbf{Repeated phase} & \textbf{Methodological focus} \\
    \hline
    Feature selection and scope clarification & Identify the next meaningful academic or administrative feature and define the intended workflow boundaries for the current increment. \\
    \hline
    Workflow and requirement analysis & Examine the target actors, conditions, information flow, constraints, and dependencies associated with the selected feature. \\
    \hline
    Design of feature behavior and user interaction & Define how the feature should behave from the user perspective, including role-sensitive interaction, states, and decision points. \\
    \hline
    Feature-slice implementation & Realize the selected feature across the relevant system layers, including interface behavior, services, permissions, and data handling where the scope requires them. \\
    \hline
    Validation and refinement & Review the implemented behavior against the intended workflow, resolve inconsistencies, and refine the feature before broader adoption. \\
    \hline
    Integration into the wider platform & Connect the completed feature to the surrounding platform so it remains coherent with existing modules, shared rules, and user workflows. \\
    \hline
  \end{tabular}
\end{table}

The iterative process was also supported by version control and pull-request
review, allowing each feature increment to be reviewed and integrated safely
before becoming part of the main codebase. In methodological terms, this
support mechanism helped keep changes bounded to one feature slice while still
checking integration effects before merge.

\section{Chapter Summary}\label{sec:chapter3_summary}

EduVerse followed an iterative, feature-based methodology with
predominantly vertical-slice development. Through this structure, the
project was developed in repeated increments centered on meaningful academic
and administrative features, with each increment progressing through
analysis, design, implementation, validation, and integration. The next
chapters move from methodology to system analysis, system design,
implementation, and testing, where the platform's structure and technical
realization are presented in greater detail.
